\documentclass[12pt, a4paper]{report}

% ─── Packages ───────────────────────────────────────────────────────────────
\usepackage[utf8]{inputenc}
\usepackage[T1]{fontenc}
\usepackage[french]{babel}
\usepackage{geometry}
\usepackage{graphicx}
\usepackage{xcolor}
\usepackage{tikz}
\usetikzlibrary{shadows, shapes.geometric}
\usepackage{titlesec}
\usepackage{fancyhdr}
\usepackage{hyperref}
\usepackage{array}
\usepackage{tabularx}
\usepackage{booktabs}
\usepackage{lmodern}
\usepackage{enumitem}
\usepackage{listings}
\usepackage{tcolorbox}
\usepackage{mdframed}
\usepackage{float}
\usepackage{caption}
\usepackage{setspace}
\usepackage{amssymb}
\usepackage{pifont}

\tcbuselibrary{skins, breakable}

\geometry{
  top=2.5cm, bottom=2.5cm,
  left=3cm, right=2.5cm
}

% ─── Colors ─────────────────────────────────────────────────────────────────
\definecolor{primaryblue}{RGB}{0, 51, 102}
\definecolor{accentred}{RGB}{180, 0, 0}
\definecolor{lightgray}{RGB}{245, 245, 245}
\definecolor{darkgray}{RGB}{80, 80, 80}
\definecolor{codebg}{RGB}{248, 248, 248}
\definecolor{owaspgreen}{RGB}{0, 120, 80}
\definecolor{critred}{RGB}{200, 0, 0}
\definecolor{highorange}{RGB}{220, 120, 0}
\definecolor{medyellow}{RGB}{180, 160, 0}
\definecolor{lowgreen}{RGB}{0, 140, 60}

% ─── Hyperref Setup ─────────────────────────────────────────────────────────
\hypersetup{
  colorlinks=true,
  linkcolor=primaryblue,
  urlcolor=accentred,
  citecolor=primaryblue,
  pdfauthor={Salma El Kamili, Fatima Ezzahrae Nouama},
  pdftitle={Scanner de Vulnérabilités Web -- OWASP Top 10},
}

% ─── Section formatting ──────────────────────────────────────────────────────
\titleformat{\chapter}[block]
  {\normalfont\LARGE\bfseries\color{primaryblue}}
  {\thechapter.}{1em}{}[{\color{accentred}\titlerule[1.5pt]}]
\titlespacing*{\chapter}{0pt}{20pt}{15pt}

\titleformat{\section}
  {\normalfont\large\bfseries\color{primaryblue}}
  {\thesection}{0.8em}{}
\titleformat{\subsection}
  {\normalfont\normalsize\bfseries\color{darkgray}}
  {\thesubsection}{0.8em}{}

% ─── Code listing style ──────────────────────────────────────────────────────
\lstdefinestyle{shellstyle}{
  backgroundcolor=\color{codebg},
  basicstyle=\ttfamily\small,
  frame=single,
  rulecolor=\color{primaryblue},
  breaklines=true,
  captionpos=b,
}
\lstset{style=shellstyle}

% ─── Custom boxes ────────────────────────────────────────────────────────────
\newtcolorbox{problembox}{
  colback=accentred!5, colframe=accentred,
  fonttitle=\bfseries, title=Problématique,
  arc=4pt, boxrule=1.5pt
}
\newtcolorbox{infobox}[1]{
  colback=primaryblue!5, colframe=primaryblue,
  fonttitle=\bfseries, title=#1,
  arc=4pt, boxrule=1pt
}

% ════════════════════════════════════════════════════════════════════════════
\begin{document}
\pagestyle{empty}

% ════════════════════════════════════════════════════════════════════════════
%  PAGE DE GARDE
% ════════════════════════════════════════════════════════════════════════════
\begin{titlepage}

  \begin{tikzpicture}[remember picture, overlay]
    \fill[primaryblue]
      (current page.north west) rectangle
      ([yshift=-4cm] current page.north east);
    \fill[primaryblue]
      (current page.south west) rectangle
      ([yshift=2cm] current page.south east);
    \fill[accentred]
      ([yshift=-4cm] current page.north west) rectangle
      ([yshift=-4.35cm] current page.north east);
    \fill[accentred]
      ([yshift=2cm] current page.south west) rectangle
      ([yshift=2.35cm] current page.south east);
  \end{tikzpicture}

  \vspace*{-1cm}
  \begin{center}
    \textcolor{white}{\large\bfseries NOM DE L'ÉTABLISSEMENT}\\[0.2cm]
    \textcolor{white}{\normalsize Filière : [Votre Filière] \quad|\quad Année académique : 2024--2025}
  \end{center}

  \vspace{0.8cm}

  \begin{center}
    \begin{tikzpicture}
      \draw[primaryblue, line width=2pt, dashed, rounded corners=5pt]
        (0,0) rectangle (4,2.5);
      \node[primaryblue] at (2,1.25) {\small Logo École};
    \end{tikzpicture}
  \end{center}

  \vspace{0.5cm}

  \begin{center}
    \begin{tikzpicture}
      \node[draw=accentred, line width=1.5pt, rounded corners=4pt,
            inner xsep=20pt, inner ysep=8pt,
            text=accentred, font=\large\bfseries] {
        Rapport de Projet de Fin d'Année (PFA)
      };
    \end{tikzpicture}
  \end{center}

  \vspace{0.6cm}

  \begin{center}
    {\color{primaryblue}\rule{0.85\linewidth}{1.5pt}}\\[0.5cm]
    {\LARGE\bfseries\color{primaryblue}
      Scanner de Vulnérabilités Web\\[0.3cm]
      \color{accentred} OWASP Top 10
    }\\[0.5cm]
    {\color{primaryblue}\rule{0.85\linewidth}{1.5pt}}
  \end{center}

  \vspace{0.5cm}

  \begin{center}
    \textcolor{darkgray}{\normalsize
      Détection automatique des vulnérabilités critiques des applications web\\
      basée sur le référentiel OWASP Top 10
    }
  \end{center}

  \vspace{0.8cm}

  \begin{center}
    \begin{tabular}{p{0.42\linewidth} p{0.42\linewidth}}
      \toprule
      \multicolumn{1}{c}{\color{primaryblue}\bfseries Réalisé par} &
      \multicolumn{1}{c}{\color{primaryblue}\bfseries Encadré par} \\
      \midrule
      \vspace{2pt}
      \textbullet\ Salma \textsc{El Kamili} \newline
      \textbullet\ Fatima Ezzahrae \textsc{Nouama} &
      \vspace{2pt}
      Pr.\ [Nom Encadrant] \\
      \bottomrule
    \end{tabular}
  \end{center}

  \vspace{0.6cm}

  \begin{center}
    \textcolor{darkgray}{\footnotesize Vulnérabilités couvertes :}\\[0.3cm]
    \footnotesize
    \begin{tabular}{@{}*{5}{>{\centering\arraybackslash}p{0.17\linewidth}}@{}}
      \colorbox{accentred!15}{\color{accentred}\bfseries A01 -- Broken Access} &
      \colorbox{accentred!15}{\color{accentred}\bfseries A02 -- Cryptographic} &
      \colorbox{accentred!15}{\color{accentred}\bfseries A03 -- Injection} &
      \colorbox{accentred!15}{\color{accentred}\bfseries A04 -- Insecure Design} &
      \colorbox{accentred!15}{\color{accentred}\bfseries A05 -- Misconfig} \\[4pt]
      \colorbox{primaryblue!12}{\color{primaryblue}\bfseries A06 -- Components} &
      \colorbox{primaryblue!12}{\color{primaryblue}\bfseries A07 -- Auth Failures} &
      \colorbox{primaryblue!12}{\color{primaryblue}\bfseries A08 -- Data Integrity} &
      \colorbox{primaryblue!12}{\color{primaryblue}\bfseries A09 -- Logging} &
      \colorbox{primaryblue!12}{\color{primaryblue}\bfseries A10 -- SSRF} \\
    \end{tabular}
  \end{center}

  \vfill
  \begin{center}
    \textcolor{white}{\normalsize Année Universitaire 2024 -- 2025}
  \end{center}
  \vspace{0.3cm}

\end{titlepage}

% ─── Header / Footer ─────────────────────────────────────────────────────────
\pagestyle{fancy}
\fancyhf{}
\fancyhead[L]{\small\color{primaryblue}\textit{Scanner de Vulnérabilités Web -- OWASP Top 10}}
\fancyhead[R]{\small\color{darkgray}\thepage}
\renewcommand{\headrulewidth}{0.4pt}
\setlength{\headheight}{14pt}

% ─── Table des matières ──────────────────────────────────────────────────────
\tableofcontents
\listoffigures
\listoftables
\newpage

% ─── Introduction générale ───────────────────────────────────────────────────
\chapter*{Introduction Générale}
\addcontentsline{toc}{chapter}{Introduction Générale}
\onehalfspacing

La sécurité des applications web représente aujourd'hui l'un des défis majeurs du
domaine informatique. Avec l'essor du numérique et la multiplication des services
en ligne, les cyberattaques ciblant les applications web se sont intensifiées,
causant des pertes financières et des atteintes à la confidentialité des données.

Dans ce contexte, ce projet de fin d'année vise à concevoir et développer un
\textbf{scanner de vulnérabilités web} automatisé, capable de détecter les failles
de sécurité les plus répandues telles que définies par le référentiel
\textbf{OWASP Top 10 2021}.

Ce rapport est structuré comme suit :
\begin{itemize}[leftmargin=2em, itemsep=3pt]
  \item Le \textbf{Chapitre 1} présente l'état de l'art : analyse des outils
        existants, des vulnérabilités et des techniques de détection.
  \item Le \textbf{Chapitre 2} expose la conception de notre solution : architecture,
        diagrammes UML et choix technologiques.
  \item Le \textbf{Chapitre 3} détaillera la réalisation et l'implémentation
        de l'outil.
\end{itemize}

\newpage

% ════════════════════════════════════════════════════════════════════════════
%  CHAPITRE 1 — ÉTAT DE L'ART
% ════════════════════════════════════════════════════════════════════════════
\chapter{État de l'Art}

\onehalfspacing

\section{Sécurité des Applications Web : Contexte Général}

La sécurité applicative est devenue un enjeu stratégique à l'ère du numérique.
Selon le rapport \textit{Verizon Data Breach Investigations Report 2023}, plus de
\textbf{74\%} des violations de données impliquent une interaction humaine ou une
exploitation de failles applicatives. Les applications web, exposées publiquement
sur Internet, constituent une surface d'attaque particulièrement vaste.

Les attaquants exploitent ces failles pour :
\begin{itemize}[leftmargin=2em, itemsep=3pt]
  \item Voler des données personnelles et financières
  \item Prendre le contrôle de comptes utilisateurs
  \item Exécuter du code arbitraire sur le serveur
  \item Accéder à des ressources internes non autorisées
\end{itemize}

\section{Le Référentiel OWASP Top 10}

\subsection{Présentation de l'OWASP}

L'\textbf{OWASP} (\textit{Open Web Application Security Project}) est une fondation
à but non lucratif qui œuvre pour améliorer la sécurité des logiciels. Elle publie
depuis 2003 la liste \textbf{OWASP Top 10}, mise à jour régulièrement pour refléter
les risques les plus critiques pour les applications web.

\subsection{OWASP Top 10 -- Édition 2021}

L'édition 2021 classe les dix catégories de risques suivantes :

\begin{center}
\begin{tabular}{>{\bfseries\centering\arraybackslash}p{1.2cm}
                >{\raggedright\arraybackslash}p{4.5cm}
                >{\raggedright\arraybackslash}p{7cm}}
  \toprule
  \textcolor{primaryblue}{Code} &
  \textcolor{primaryblue}{Catégorie} &
  \textcolor{primaryblue}{Description} \\
  \midrule
  A01 & Broken Access Control
      & Contrôle d'accès insuffisant permettant des accès non autorisés aux ressources \\[3pt]
  A02 & Cryptographic Failures
      & Mauvaise utilisation ou absence de chiffrement exposant des données sensibles \\[3pt]
  A03 & Injection
      & Injection de code malveillant (SQL, NoSQL, OS, LDAP) via des entrées non validées \\[3pt]
  A04 & Insecure Design
      & Failles architecturales dues à l'absence de modélisation des menaces \\[3pt]
  A05 & Security Misconfiguration
      & Configurations par défaut non sécurisées, services inutiles activés \\[3pt]
  A06 & Vulnerable Components
      & Utilisation de bibliothèques et frameworks avec des failles connues \\[3pt]
  A07 & Auth \& Session Failures
      & Authentification faible, gestion de session déficiente \\[3pt]
  A08 & Software \& Data Integrity
      & Absence de vérification d'intégrité sur les mises à jour et pipelines CI/CD \\[3pt]
  A09 & Logging \& Monitoring Failures
      & Journalisation insuffisante retardant la détection des attaques \\[3pt]
  A10 & SSRF
      & Falsification de requêtes côté serveur vers des ressources internes \\
  \bottomrule
\end{tabular}
\captionof{table}{OWASP Top 10 -- Édition 2021}
\end{center}

\section{Vulnérabilités Ciblées : Définitions et Mécanismes}

\subsection{SQL Injection (SQLi)}

L'injection SQL est l'une des vulnérabilités les plus anciennes et les plus
dangereuses. Elle survient lorsque des entrées utilisateur non filtrées sont
incorporées directement dans des requêtes SQL. L'attaquant peut alors manipuler
la requête pour extraire, modifier ou supprimer des données.

\begin{lstlisting}[language=SQL, caption={Exemple de payload SQLi classique}]
' OR '1'='1
' UNION SELECT username, password FROM users --
\end{lstlisting}

\subsection{Cross-Site Scripting (XSS)}

Le XSS permet à un attaquant d'injecter du code JavaScript malveillant dans
les pages vues par d'autres utilisateurs. On distingue :
\begin{itemize}[leftmargin=2em, itemsep=2pt]
  \item \textbf{XSS Réfléchi :} le payload est inclus dans la réponse HTTP
        immédiate (via un paramètre d'URL).
  \item \textbf{XSS Persistant :} le payload est stocké en base de données
        et servi à tous les visiteurs.
  \item \textbf{XSS basé sur le DOM :} manipulation directe du DOM côté client.
\end{itemize}

\subsection{IDOR (Insecure Direct Object Reference)}

L'IDOR se produit lorsqu'une application expose des références directes à des
objets internes (identifiants, noms de fichiers) sans vérifier les autorisations
de l'utilisateur. Un attaquant peut modifier ces références pour accéder à des
ressources appartenant à d'autres utilisateurs.

\subsection{SSRF (Server-Side Request Forgery)}

Le SSRF force le serveur à effectuer des requêtes HTTP vers des ressources
internes ou externes choisies par l'attaquant. Cela peut permettre d'atteindre
des services internes normalement inaccessibles depuis l'extérieur (métadonnées
cloud, bases de données internes, etc.).

\subsection{Path Traversal}

Cette attaque exploite une validation insuffisante des chemins de fichiers pour
accéder à des fichiers arbitraires sur le serveur (ex : \texttt{../../etc/passwd}).

\subsection{XXE (XML External Entity)}

Lorsqu'un parseur XML traite des entités externes définies par l'attaquant,
il peut divulguer des fichiers locaux, effectuer des SSRF, ou provoquer un déni
de service.

\section{Outils de Scan de Vulnérabilités Existants}

\subsection{OWASP ZAP (Zed Attack Proxy)}

OWASP ZAP est un outil open source développé par l'OWASP, disponible en mode
proxy intercepteur ou en ligne de commande. Il propose un scan actif et passif,
et dispose d'une API REST permettant l'intégration dans des pipelines CI/CD.

\textbf{Points forts :} open source, communauté active, extensible via plugins.\\
\textbf{Limites :} interface complexe, lent sur les grandes applications.

\subsection{Burp Suite}

Burp Suite (PortSwigger) est la référence professionnelle pour les tests
d'intrusion web. La version Community est gratuite mais limitée ; la version
Professional offre un scanner automatisé puissant.

\textbf{Points forts :} très complet, fiable, outils manuels avancés.\\
\textbf{Limites :} version complète payante (coût élevé), fermé.

\subsection{Nikto}

Nikto est un scanner web léger en ligne de commande, orienté détection de
configurations dangereuses, versions obsolètes et fichiers sensibles exposés.

\textbf{Points forts :} rapide, simple, open source.\\
\textbf{Limites :} peu de couverture des vulnérabilités applicatives (SQLi, XSS...).

\subsection{Acunetix / Invicti}

Acunetix est un scanner commercial performant avec une couverture étendue de
l'OWASP Top 10, une interface web et des rapports détaillés.

\textbf{Points forts :} couverture étendue, rapports professionnels.\\
\textbf{Limites :} prix très élevé, non accessible pour les projets académiques.

\subsection{Tableau Comparatif des Outils}

\begin{center}
\begin{tabular}{l c c c c c}
  \toprule
  \textbf{Outil} & \textbf{Open Source} & \textbf{SQLi} & \textbf{XSS}
                 & \textbf{SSRF} & \textbf{CLI} \\
  \midrule
  OWASP ZAP    & \textcolor{lowgreen}{$\checkmark$} & \textcolor{lowgreen}{$\checkmark$}
               & \textcolor{lowgreen}{$\checkmark$} & Partiel
               & \textcolor{lowgreen}{$\checkmark$} \\
  Burp Suite   & \textcolor{accentred}{$\times$}  & \textcolor{lowgreen}{$\checkmark$}
               & \textcolor{lowgreen}{$\checkmark$} & \textcolor{lowgreen}{$\checkmark$}
               & Partiel \\
  Nikto        & \textcolor{lowgreen}{$\checkmark$} & \textcolor{accentred}{$\times$}
               & \textcolor{accentred}{$\times$}  & \textcolor{accentred}{$\times$}
               & \textcolor{lowgreen}{$\checkmark$} \\
  Acunetix     & \textcolor{accentred}{$\times$}  & \textcolor{lowgreen}{$\checkmark$}
               & \textcolor{lowgreen}{$\checkmark$} & \textcolor{lowgreen}{$\checkmark$}
               & \textcolor{accentred}{$\times$}  \\
  \textbf{Notre outil} & \textcolor{lowgreen}{$\checkmark$} & \textcolor{lowgreen}{$\checkmark$}
               & \textcolor{lowgreen}{$\checkmark$} & \textcolor{lowgreen}{$\checkmark$}
               & \textcolor{lowgreen}{$\checkmark$} \\
  \bottomrule
\end{tabular}
\captionof{table}{Comparaison des outils de scan de vulnérabilités web}
\end{center}

\section{Techniques de Détection Automatique}

\subsection{Scan Actif vs Scan Passif}

\begin{itemize}[leftmargin=2em, itemsep=4pt]
  \item \textbf{Scan passif :} analyse le trafic HTTP sans envoyer de requêtes
        supplémentaires. Non intrusif mais limité en couverture.
  \item \textbf{Scan actif :} envoie des requêtes spécialement forgées pour
        provoquer des comportements révélateurs. Plus efficace mais plus intrusif.
\end{itemize}

Notre outil adopte une approche de \textbf{scan actif ciblé}.

\subsection{Fuzzing}

Le fuzzing consiste à envoyer des données semi-aléatoires ou basées sur des
dictionnaires (\textit{wordlists}) dans les paramètres d'entrée afin de provoquer
des erreurs ou des comportements anormaux révélant des vulnérabilités.

\subsection{Analyse par Signatures}

Chaque type de vulnérabilité possède des signatures caractéristiques dans les
réponses HTTP. Par exemple, une erreur SQL de type \texttt{syntax error near...}
indique une potentielle injection SQL. Notre moteur analyse ces signatures pour
confirmer la présence d'une faille.

\section{Synthèse et Positionnement}

Les outils existants sont soit trop génériques (Nikto), soit trop coûteux
(Burp Suite Pro, Acunetix), soit complexes à intégrer dans un workflow étudiant.
Notre projet se positionne comme un outil \textbf{open source, léger, modulaire
et pédagogique}, focalisé sur les vulnérabilités de l'OWASP Top 10 et utilisable
directement en ligne de commande avec Python.

\begin{infobox}{Positionnement de notre outil}
\begin{itemize}[leftmargin=1.5em, itemsep=2pt]
  \item Gratuit et open source
  \item Orienté OWASP Top 10 (détection ciblée et précise)
  \item Interface CLI simple : \texttt{python scanner.py --url <cible>}
  \item Rapport HTML exploitable immédiatement
  \item Architecture modulaire facilitant les extensions futures
\end{itemize}
\end{infobox}

\newpage

% ════════════════════════════════════════════════════════════════════════════
%  CHAPITRE 2 — CONCEPTION
% ════════════════════════════════════════════════════════════════════════════
\chapter{Conception}

\onehalfspacing

\section{Architecture Générale du Système}

Notre scanner est organisé selon une \textbf{architecture en pipeline à quatre modules}
indépendants et communicants. Cette approche modulaire facilite la maintenance,
les tests unitaires et l'ajout de nouveaux modules de détection.

\begin{figure}[H]
\centering
\begin{tikzpicture}[
  box/.style={draw=primaryblue, fill=primaryblue!8, rounded corners=6pt,
              minimum width=2.8cm, minimum height=1.4cm,
              font=\small\bfseries\color{primaryblue}, align=center,
              drop shadow},
  arr/.style={->, very thick, color=accentred},
  lbl/.style={font=\tiny\color{darkgray}, above=2pt}
]
  \node[box] (input)  at (0,0)    {Entrée\\{\tiny URL cible}};
  \node[box] (crawl)  at (3.5,0)  {Crawling\\{\tiny \& Discovery}};
  \node[box] (scan)   at (7,0)    {Scan\\{\tiny \& Détection}};
  \node[box] (eval)   at (10.5,0) {Évaluation\\{\tiny CVSS v3.1}};
  \node[box] (report) at (7,-2.5) {Reporting\\{\tiny HTML}};

  \draw[arr] (input)  -- node[lbl]{URL} (crawl);
  \draw[arr] (crawl)  -- node[lbl]{Endpoints} (scan);
  \draw[arr] (scan)   -- node[lbl]{Vulnérabilités} (eval);
  \draw[arr] (eval)   -- node[lbl]{Scores} ++(0,-1.2) -| (report);
\end{tikzpicture}
\caption{Architecture en pipeline du scanner de vulnérabilités}
\end{figure}

\section{Diagramme de Cas d'Utilisation}

\begin{figure}[H]
\centering
\begin{tikzpicture}[
  uc/.style={draw=primaryblue, fill=primaryblue!6, rounded corners=4pt,
             minimum width=4.5cm, minimum height=0.7cm,
             font=\small, align=center},
  actor/.style={font=\small\bfseries\color{primaryblue}},
  arr/.style={->, thick, color=darkgray}
]
  % Actor
  \draw[primaryblue, line width=1.5pt] (0,0) circle (0.3cm);
  \draw[primaryblue, line width=1.5pt] (0,-0.3) -- (0,-1);
  \draw[primaryblue, line width=1.5pt] (-0.4,-0.6) -- (0.4,-0.6);
  \draw[primaryblue, line width=1.5pt] (0,-1) -- (-0.3,-1.5);
  \draw[primaryblue, line width=1.5pt] (0,-1) -- (0.3,-1.5);
  \node[actor] at (0,-2) {Utilisateur};

  % Use cases
  \node[uc] (uc1) at (5, 1.8)  {Lancer un scan};
  \node[uc] (uc2) at (5, 0.6)  {Configurer les options};
  \node[uc] (uc3) at (5,-0.6)  {Consulter les résultats};
  \node[uc] (uc4) at (5,-1.8)  {Exporter le rapport HTML};

  % Arrows
  \draw[arr] (0.4,-0.75) -- (uc1.west);
  \draw[arr] (0.4,-0.75) -- (uc2.west);
  \draw[arr] (0.4,-0.75) -- (uc3.west);
  \draw[arr] (0.4,-0.75) -- (uc4.west);

  % System boundary
  \draw[primaryblue, dashed, rounded corners=6pt, line width=1pt]
    (3.2, 2.4) rectangle (9.3,-2.4);
  \node[primaryblue, font=\small\bfseries] at (6.25, 2.1) {Système Scanner};
\end{tikzpicture}
\caption{Diagramme de cas d'utilisation}
\end{figure}

\section{Diagramme de Classes}

\begin{figure}[H]
\centering
\begin{tikzpicture}[
  cls/.style={draw=primaryblue, fill=white, rounded corners=3pt,
              minimum width=3.8cm, font=\small},
  arr/.style={->, thick, color=darkgray},
  comp/.style={->, thick, color=accentred, dashed}
]
  % Scanner (main)
  \node[cls] (scanner) at (5.5,6) {
    \begin{tabular}{|l|}
      \hline
      \multicolumn{1}{|c|}{\textbf{\color{primaryblue}Scanner}} \\
      \hline
      url : str \\
      options : dict \\
      \hline
      run() \\
      get\_report() \\
      \hline
    \end{tabular}
  };

  % Crawler
  \node[cls] (crawler) at (0,3) {
    \begin{tabular}{|l|}
      \hline
      \multicolumn{1}{|c|}{\textbf{\color{primaryblue}Crawler}} \\
      \hline
      base\_url : str \\
      visited : set \\
      \hline
      crawl() \\
      extract\_forms() \\
      \hline
    \end{tabular}
  };

  % VulnDetector
  \node[cls] (detector) at (5.5,3) {
    \begin{tabular}{|l|}
      \hline
      \multicolumn{1}{|c|}{\textbf{\color{primaryblue}VulnDetector}} \\
      \hline
      payloads : list \\
      \hline
      detect\_sqli() \\
      detect\_xss() \\
      detect\_ssrf() \\
      detect\_idor() \\
      detect\_xxe() \\
      \hline
    \end{tabular}
  };

  % Evaluator
  \node[cls] (eval) at (11,3) {
    \begin{tabular}{|l|}
      \hline
      \multicolumn{1}{|c|}{\textbf{\color{primaryblue}CVSSEvaluator}} \\
      \hline
      base\_score : float \\
      \hline
      compute\_score() \\
      get\_severity() \\
      \hline
    \end{tabular}
  };

  % Reporter
  \node[cls] (reporter) at (5.5,0) {
    \begin{tabular}{|l|}
      \hline
      \multicolumn{1}{|c|}{\textbf{\color{primaryblue}HTMLReporter}} \\
      \hline
      vulns : list \\
      output\_path : str \\
      \hline
      generate() \\
      save() \\
      \hline
    \end{tabular}
  };

  % Vulnerability
  \node[cls] (vuln) at (11,0) {
    \begin{tabular}{|l|}
      \hline
      \multicolumn{1}{|c|}{\textbf{\color{primaryblue}Vulnerability}} \\
      \hline
      name : str \\
      url : str \\
      payload : str \\
      severity : str \\
      cvss : float \\
      \hline
      to\_dict() \\
      \hline
    \end{tabular}
  };

  % Relations
  \draw[comp] (scanner) -- (crawler);
  \draw[comp] (scanner) -- (detector);
  \draw[comp] (scanner) -- (eval);
  \draw[comp] (scanner) -- (reporter);
  \draw[arr]  (detector) -- node[right,font=\tiny]{crée} (vuln);
  \draw[arr]  (eval)     -- (vuln);
  \draw[arr]  (reporter) -- (vuln);
\end{tikzpicture}
\caption{Diagramme de classes du scanner}
\end{figure}

\section{Diagramme de Séquence -- Scan d'une Vulnérabilité}

\begin{figure}[H]
\centering
\begin{tikzpicture}[
  lifeline/.style={draw=primaryblue, dashed, very thin},
  msg/.style={->, thick, color=primaryblue},
  ret/.style={->, dashed, thick, color=darkgray},
  box/.style={draw=primaryblue, fill=primaryblue!10, rounded corners=2pt,
              font=\small\bfseries\color{primaryblue}, minimum width=2cm,
              minimum height=0.5cm, align=center}
]
  % Lifeline headers
  \node[box] (user)     at (0,0)    {Utilisateur};
  \node[box] (scanner)  at (3,0)    {Scanner};
  \node[box] (crawler)  at (6,0)    {Crawler};
  \node[box] (detector) at (9,0)    {Detector};
  \node[box] (reporter) at (12,0)   {Reporter};

  % Lifelines
  \foreach \x in {0,3,6,9,12}
    \draw[lifeline] (\x,-0.3) -- (\x,-7);

  % Messages
  \draw[msg] (0,-0.8)  -- node[above,font=\tiny]{run(url)} (3,-0.8);
  \draw[msg] (3,-1.4)  -- node[above,font=\tiny]{crawl(url)} (6,-1.4);
  \draw[ret] (6,-2.0)  -- node[above,font=\tiny]{endpoints, forms} (3,-2.0);
  \draw[msg] (3,-2.6)  -- node[above,font=\tiny]{detect(endpoints)} (9,-2.6);
  \draw[ret] (9,-3.2)  -- node[above,font=\tiny]{vulns + payloads} (3,-3.2);
  \draw[msg] (3,-3.8)  -- node[above,font=\tiny]{evaluate(vulns)} (9,-3.8);
  \draw[ret] (9,-4.4)  -- node[above,font=\tiny]{scored vulns} (3,-4.4);
  \draw[msg] (3,-5.0)  -- node[above,font=\tiny]{generate\_report()} (12,-5.0);
  \draw[ret] (12,-5.6) -- node[above,font=\tiny]{rapport.html} (3,-5.6);
  \draw[ret] (3,-6.2)  -- node[above,font=\tiny]{rapport.html} (0,-6.2);
\end{tikzpicture}
\caption{Diagramme de séquence d'un cycle de scan complet}
\end{figure}

\section{Description Détaillée des Modules}

\subsection{Module de Crawling}

Le module de crawling explore récursivement l'application cible à partir de l'URL
de base fournie. Il utilise la bibliothèque \texttt{requests} pour les requêtes HTTP
et \texttt{BeautifulSoup} pour l'analyse du HTML.

\textbf{Fonctionnement :}
\begin{enumerate}[leftmargin=2em, itemsep=3pt]
  \item Téléchargement de la page à l'URL de départ
  \item Extraction de tous les liens \texttt{<a href>} internes au domaine
  \item Identification des formulaires \texttt{<form>} et de leurs champs
  \item Collecte des paramètres GET depuis les URLs
  \item Mise en file d'attente des nouvelles URLs à visiter (BFS)
\end{enumerate}

\subsection{Module de Détection}

Chaque type de vulnérabilité est géré par une méthode dédiée. Le module charge
des listes de \textit{payloads} depuis des fichiers de configuration et les injecte
dans chaque paramètre identifié.

\begin{center}
\begin{tabular}{l >{\raggedright\arraybackslash}p{5cm} >{\raggedright\arraybackslash}p{4cm}}
  \toprule
  \textbf{Vulnérabilité} & \textbf{Méthode de détection} & \textbf{Indicateur de succès} \\
  \midrule
  SQL Injection & Payload avec guillemets, UNION, sleep & Message d'erreur SQL, délai \\
  XSS           & Injection \texttt{<script>alert()</script>} & Réflexion dans la réponse \\
  IDOR          & Incrémentation d'ID dans les paramètres & Contenu d'un autre utilisateur \\
  SSRF          & URL pointant vers 127.0.0.1 ou metadata & Réponse interne reçue \\
  Path Traversal& \texttt{../../etc/passwd} dans les paramètres fichier & Contenu système dans réponse \\
  XXE           & Entité XML externe dans les corps XML & Contenu de fichier local \\
  \bottomrule
\end{tabular}
\captionof{table}{Méthodes de détection par type de vulnérabilité}
\end{center}

\subsection{Module d'Évaluation (CVSS v3.1)}

Chaque vulnérabilité détectée reçoit un score CVSS calculé selon les métriques
de base : vecteur d'attaque, complexité, privilèges requis, interaction utilisateur,
et impact sur la confidentialité, l'intégrité et la disponibilité.

\begin{center}
\begin{tabular}{l c l}
  \toprule
  \textbf{Score CVSS} & \textbf{Plage} & \textbf{Niveau} \\
  \midrule
  \textcolor{critred}{\textbf{Critique}}    & 9.0 -- 10.0 & Correction immédiate requise \\
  \textcolor{highorange}{\textbf{Élevé}}    & 7.0 -- 8.9  & Correction prioritaire \\
  \textcolor{medyellow}{\textbf{Moyen}}     & 4.0 -- 6.9  & Correction planifiée \\
  \textcolor{lowgreen}{\textbf{Faible}}     & 0.1 -- 3.9  & Correction à terme \\
  \bottomrule
\end{tabular}
\captionof{table}{Classification de gravité selon CVSS v3.1}
\end{center}

\subsection{Module de Reporting HTML}

Le module génère un fichier HTML autonome contenant :
\begin{itemize}[leftmargin=2em, itemsep=3pt]
  \item Un tableau de bord avec le nombre de vulnérabilités par niveau de gravité
  \item Une liste détaillée de chaque faille avec URL, payload, preuve et recommandation
  \item Un style CSS intégré pour un rendu lisible sans dépendance externe
\end{itemize}

\section{Technologies et Environnement Technique}

\begin{center}
\begin{tabular}{l l >{\raggedright\arraybackslash}p{6cm}}
  \toprule
  \textbf{Composant} & \textbf{Technologie} & \textbf{Rôle} \\
  \midrule
  Langage           & Python 3.10+        & Développement de l'outil \\
  Requêtes HTTP     & \texttt{requests}   & Envoi et réception des requêtes web \\
  Parsing HTML      & \texttt{BeautifulSoup4} & Extraction des liens et formulaires \\
  Analyse réponses  & \texttt{re} (regex) & Détection de signatures de vulnérabilités \\
  Rapport           & HTML/CSS pur        & Génération du rapport final \\
  CLI               & \texttt{argparse}   & Interface ligne de commande \\
  Tests             & \texttt{pytest}     & Tests unitaires des modules \\
  Cible de test     & DVWA / WebGoat      & Applications vulnérables intentionnellement \\
  \bottomrule
\end{tabular}
\captionof{table}{Stack technologique du projet}
\end{center}

\section{Structure du Projet}

\begin{lstlisting}[caption={Arborescence du projet}]
scanner/
|-- scanner.py          # Point d'entree CLI
|-- crawler/
|   |-- __init__.py
|   `-- crawler.py      # Module de crawling
|-- detectors/
|   |-- __init__.py
|   |-- sqli.py         # Detecteur SQLi
|   |-- xss.py          # Detecteur XSS
|   |-- ssrf.py         # Detecteur SSRF
|   |-- idor.py         # Detecteur IDOR
|   |-- xxe.py          # Detecteur XXE
|   `-- path_traversal.py
|-- evaluator/
|   `-- cvss.py         # Calcul du score CVSS
|-- reporter/
|   `-- html_reporter.py
|-- payloads/
|   |-- sqli.txt
|   |-- xss.txt
|   `-- path_traversal.txt
`-- tests/
    `-- test_detectors.py
\end{lstlisting}

\newpage

% ════════════════════════════════════════════════════════════════════════════
%  CONCLUSION GÉNÉRALE
% ════════════════════════════════════════════════════════════════════════════
\chapter*{Conclusion Générale}
\addcontentsline{toc}{chapter}{Conclusion Générale}

\onehalfspacing

Ce rapport a présenté les fondements théoriques et la conception de notre scanner
de vulnérabilités web basé sur l'OWASP Top 10. Le premier chapitre a dressé un
état de l'art approfondi, comparant les outils existants et définissant les
mécanismes des vulnérabilités ciblées. Le second chapitre a formalisé notre
conception via des diagrammes UML (cas d'utilisation, classes, séquence) et
précisé l'architecture modulaire retenue.

La phase suivante, objet du chapitre 3, consistera en la \textbf{réalisation et
l'implémentation} de l'outil : développement des modules de crawling, de détection,
d'évaluation CVSS et de génération de rapports HTML, ainsi que les tests sur des
environnements volontairement vulnérables (DVWA, WebGoat).

% ─── Bibliographie ───────────────────────────────────────────────────────────
\begin{thebibliography}{9}

\bibitem{owasp2021}
OWASP Foundation,
\textit{OWASP Top 10 -- 2021},
\url{https://owasp.org/Top10/}, 2021.

\bibitem{verizon2023}
Verizon,
\textit{2023 Data Breach Investigations Report},
\url{https://www.verizon.com/business/resources/reports/dbir/}, 2023.

\bibitem{cvss31}
FIRST,
\textit{Common Vulnerability Scoring System v3.1: Specification Document},
\url{https://www.first.org/cvss/v3.1/specification-document}, 2019.

\bibitem{zap}
OWASP Foundation,
\textit{OWASP ZAP -- Zed Attack Proxy},
\url{https://www.zaproxy.org/}, 2023.

\bibitem{burp}
PortSwigger,
\textit{Burp Suite Web Vulnerability Scanner},
\url{https://portswigger.net/burp}, 2023.

\bibitem{nikto}
Nikto Project,
\textit{Nikto Web Server Scanner},
\url{https://cirt.net/Nikto2}, 2023.

\bibitem{requests}
K. Reitz,
\textit{Requests: HTTP for Humans},
\url{https://docs.python-requests.org/}, 2023.

\bibitem{bs4}
L. Richardson,
\textit{Beautiful Soup Documentation},
\url{https://www.crummy.com/software/BeautifulSoup/}, 2023.

\end{thebibliography}

\end{document}